% !TeX root = ../main.tex
\chapter{Applicazioni lineari}
\label{cap:applicazioni-lineari}
\fonti{Lezioni 16, 17 e 18}

Una matrice permette di calcolare combinazioni lineari delle sue colonne. Le applicazioni lineari descrivono questa stessa struttura fra spazi vettoriali anche quando i vettori non sono già presentati come colonne. Le basi consentono di passare dall'applicazione astratta al calcolo matriciale.
% !TeX root = ../main.tex
\section{Linearità, composizione e isomorfismi}
\begin{definizione}
Siano $V,W$ spazi vettoriali sullo stesso campo $K$. Una funzione $f:V\to W$ è \emph{lineare}, o un \emph{omomorfismo di spazi vettoriali}, se
\[
 f(v_1+v_2)=f(v_1)+f(v_2),\qquad f(\lambda v)=\lambda f(v)
\]
per ogni $v_1,v_2,v\in V$ e $\lambda\in K$. Equivalentemente, per tutti $v_1,v_2\in V$ e $\lambda_1,\lambda_2\in K$,
\[
 f(\lambda_1v_1+\lambda_2v_2)=\lambda_1f(v_1)+\lambda_2f(v_2).
\]
\end{definizione}
L'equivalenza segue applicando prima l'additività e poi l'omogeneità; viceversa, basta scegliere prima $\lambda_1=\lambda_2=1$ e poi $\lambda_2=0$. Iterando, una funzione lineare conserva ogni combinazione lineare finita.

\begin{osservazione}
Ogni applicazione lineare soddisfa $f(0_V)=0_W$: da $f(0_V)=f(0_V+0_V)=f(0_V)+f(0_V)$ basta sottrarre $f(0_V)$. Inoltre $f(-v)=-f(v)$ e $f(v_1-v_2)=f(v_1)-f(v_2)$. La condizione $f(0_V)=0_W$ è necessaria, ma non sufficiente per la linearità.
\end{osservazione}
\begin{esempio}
Le seguenti funzioni da $\R$ a $\R$ non sono lineari.
\begin{itemize}
\item $f(x)=x+1$: si ha $f(0)=1$.
\item $f(x)=\sin x$: pur essendo $f(0)=0$, scegliendo $x_1=x_2=\pi/2$ si trova $f(x_1+x_2)=0\ne2=f(x_1)+f(x_2)$.
\item $f(x)=x^2$: si ha $f(1+1)=4\ne2=f(1)+f(1)$. Più in generale $x\mapsto x^k$, per un intero $k\geq2$, non è lineare su $\R$, perché $2^k\ne2$.
\end{itemize}
In coordinate reali le applicazioni lineari hanno componenti date da espressioni omogenee di primo grado, senza termine costante; una componente può anche essere identicamente nulla.\footnote{La forma scritta di una funzione polinomiale può ingannare su campi finiti: sul campo $\mathbf F_2=\{0,1\}$ la funzione $x\mapsto x^2$ coincide con l'identità, perché $0^2=0$ e $1^2=1$, ed è quindi lineare.}
\end{esempio}

\begin{proposizione}
La composizione di applicazioni lineari è lineare. Se un'applicazione lineare $f:V\to W$ è biiettiva, anche la sua inversa $f^{-1}:W\to V$ è lineare.
\end{proposizione}
\begin{proof}
\passomath{g\circ f} Per $f:V\to W$ e $g:W\to Z$ lineari,
\begin{align*}
 (g\circ f)(\lambda v+\mu v')
 &=g\bigl(\lambda f(v)+\mu f(v')\bigr)\\
 &=\lambda(g\circ f)(v)+\mu(g\circ f)(v').
\end{align*}
\passomath{f^{-1}} Siano $w_1,w_2\in W$ e poniamo $v_i=f^{-1}(w_i)$. Per linearità $f(\lambda v_1+\mu v_2)=\lambda w_1+\mu w_2$. Applicando l'inversa otteniamo
\[
 f^{-1}(\lambda w_1+\mu w_2)=\lambda f^{-1}(w_1)+\mu f^{-1}(w_2).
\]
\end{proof}
\begin{definizione}
Un'applicazione lineare biiettiva si chiama \emph{isomorfismo}. Due spazi si dicono \emph{isomorfi} se esiste un isomorfismo fra essi.
\end{definizione}

\section{Nucleo, immagine e iniettività}
\begin{definizione}
Per $f:V\to W$ lineare, il \emph{nucleo} e l'\emph{immagine} sono
\[
 \ker f=\{v\in V:f(v)=0_W\},\qquad
 \operatorname{Im}f=\{f(v):v\in V\}\subseteq W.
\]
Il nucleo si indica anche con $N(f)$. L'applicazione è suriettiva se e solo se $\operatorname{Im}f=W$.
\end{definizione}
\begin{teorema}
Il nucleo è un sottospazio di $V$ e l'immagine è un sottospazio di $W$.
\end{teorema}
\begin{proof}
\passomath{\ker f} Il vettore nullo appartiene al nucleo. Se $v_1,v_2\in\ker f$, allora $f(\lambda v_1+\mu v_2)=\lambda0_W+\mu0_W=0_W$: il nucleo è chiuso per combinazioni lineari.

\passomath{\operatorname{Im}f} L'immagine contiene $0_W=f(0_V)$. Se $w_i=f(v_i)$, allora $\lambda w_1+\mu w_2=f(\lambda v_1+\mu v_2)$ appartiene ancora all'immagine. Vale quindi il criterio di sottospazio anche per $\operatorname{Im}f$.
\end{proof}
\begin{teorema}
Un'applicazione lineare $f:V\to W$ è iniettiva se e solo se $\ker f=\{0_V\}$.
\end{teorema}
\begin{proof}
\passomath{\Rightarrow} Se $v\in\ker f$, allora $f(v)=0_W=f(0_V)$; l'iniettività dà $v=0_V$.

\passomath{\Leftarrow} Se $f(v_1)=f(v_2)$, per linearità $f(v_1-v_2)=0_W$. Dunque $v_1-v_2\in\ker f=\{0_V\}$, ossia $v_1=v_2$.
\end{proof}

\section{Coordinate e teorema della dimensione}
\subsection{L'isomorfismo delle coordinate}
\begin{esempio}
Lo spazio $\R[X]_{\leq3}$ dei polinomi di grado al più $3$ è isomorfo a $\R^4$ mediante
\[
 a_0+a_1X+a_2X^2+a_3X^3\longmapsto(a_0,a_1,a_2,a_3)^T.
\]
La somma di polinomi e la moltiplicazione per scalari agiscono sui coefficienti componente per componente, quindi la funzione è lineare. È biiettiva perché ogni quadrupla determina uno e un solo polinomio. Questo formalizza il passaggio alle coordinate già usato per le equazioni dei sottospazi di polinomi.
\end{esempio}
\begin{teorema}
Se $V$ ha dimensione finita $n$ su $K$ e $\mathcal B=(v_1,\ldots,v_n)$ è una base ordinata, la funzione
\[
 c_{\mathcal B}:V\longrightarrow K^n,\qquad v\longmapsto[v]_{\mathcal B}
\]
è un isomorfismo. In particolare $V$ è isomorfo a $K^n$.
\end{teorema}
\begin{proof}
Se $v=\sum_i\lambda_iv_i$ e $v'=\sum_i\mu_iv_i$, allora $av+bv'=\sum_i(a\lambda_i+b\mu_i)v_i$. Per unicità delle coordinate,
\[
 [av+bv']_{\mathcal B}=a[v]_{\mathcal B}+b[v']_{\mathcal B}.
\]
La mappa è iniettiva perché coordinate uguali ricostruiscono lo stesso vettore, ed è suriettiva perché ogni colonna $(\lambda_i)$ è immagine di $\sum_i\lambda_iv_i$.
\end{proof}
L'isomorfismo dipende dalla base scelta: non è canonico in assenza di una scelta privilegiata di base.

\subsection{Rango e nullità di un'applicazione}
\begin{teorema}[Teorema della dimensione]
\label{teo:dimensione-applicazione}
Sia $f:V\to W$ lineare con $\dim V=n<\infty$. Allora
\[
 \dim\ker f+\dim\operatorname{Im}f=n.
\]
\end{teorema}
\begin{proof}
\textbf{Scelta della base.} Sia $(v_1,\ldots,v_r)$ una base di $\ker f$. Completiamola a una base $(v_1,\ldots,v_r,v_{r+1},\ldots,v_n)$ di $V$.

\textbf{Generazione dell'immagine.} Se $w\in\operatorname{Im}f$, esiste $v=\sum_{i=1}^n\lambda_iv_i$ tale che $w=f(v)$. Le prime $r$ immagini sono nulle, quindi
\[
 w=\sum_{i=r+1}^n\lambda_if(v_i).
\]
Le immagini $f(v_{r+1}),\ldots,f(v_n)$ generano dunque $\operatorname{Im}f$.

\textbf{Indipendenza.} Supponiamo $\sum_{i=r+1}^n\alpha_if(v_i)=0_W$. Per linearità $\sum_{i=r+1}^n\alpha_iv_i\in\ker f$, quindi si può scrivere come $\sum_{j=1}^r\beta_jv_j$. Ne segue
\[
 \sum_{i=r+1}^n\alpha_iv_i-\sum_{j=1}^r\beta_jv_j=0_V.
\]
Poiché $(v_1,\ldots,v_n)$ è una base, tutti i coefficienti sono nulli. Le $n-r$ immagini sono pertanto una base dell'immagine e $\dim\operatorname{Im}f=n-r$. Il ragionamento comprende i casi $r=0$ e $r=n$, con le famiglie vuote opportune.
\end{proof}
\begin{definizione}
La \emph{nullità} di $f$ è $\dim\ker f$; il \emph{rango} di $f$ è $\dim\operatorname{Im}f$. Il teorema afferma che nullità e rango sommano alla dimensione del dominio.
\end{definizione}

% !TeX root = ../main.tex
\section{Proiezioni e simmetrie rispetto a una somma diretta}
Sia $V=U\oplus W$. La somma diretta assicura che ogni $v\in V$ si scriva in modo unico come $v=u+w$, con $u\in U$ e $w\in W$. Questa unicità permette di definire applicazioni a partire dalle due componenti.

\subsection{Proiezione su un sottospazio lungo un complementare}
\begin{definizione}
La \emph{proiezione su $U$ nella direzione di $W$} è
\[
 \pi_U^W:V\longrightarrow V,\qquad \pi_U^W(u+w)=u.
\]
\end{definizione}
\begin{schema}
\begin{tikzpicture}[scale=.9]
 \coordinate (O) at (0,0); \coordinate (U) at (3,0);
 \coordinate (W) at (1,1.6); \coordinate (V) at (4,1.6);
 \draw[asse] (-.6,0)--(4.6,0) node[right] {$U$};
 \draw[asse] (-.3,-.48)--(1.5,2.4) node[above] {$W$};
 \draw[dashed,black!45] (U)--(V)--(W);
 \draw[vettore] (O)--(V) node[right] {$v=u+w$};
 \draw[vettore] (O)--(U) node[below] {$u=\pi_U^W(v)$};
 \draw[vettore] (O)--(W) node[left] {$w$};
\end{tikzpicture}
\end{schema}
Non è richiesta ortogonalità fra $U$ e $W$: la direzione è fissata dal complementare scelto.
\begin{proposizione}
La proiezione $\pi=\pi_U^W$ è lineare e soddisfa
\[
 \ker\pi=W,\qquad \operatorname{Im}\pi=U,\qquad
 \pi(u)=u\ (u\in U),\qquad \pi\circ\pi=\pi.
\]
\end{proposizione}
\begin{proof}
\textbf{Linearità.} Per $v_i=u_i+w_i$, la scomposizione di $\lambda v_1+\mu v_2$ è $(\lambda u_1+\mu u_2)+(\lambda w_1+\mu w_2)$. Dunque
\[
 \pi(\lambda v_1+\mu v_2)=\lambda u_1+\mu u_2=\lambda\pi(v_1)+\mu\pi(v_2).
\]
\textbf{Nucleo e immagine.} Si ha $\pi(v)=0$ se e solo se $u=0$, ossia $v\in W$. Ogni immagine appartiene a $U$ e ogni $u\in U$ è immagine di se stesso, perché $u=u+0_W$.

\textbf{Idempotenza.} Poiché $\pi(v)\in U$ e la proiezione fissa ogni vettore di $U$, si ha $\pi(\pi(v))=\pi(v)$.
\end{proof}

\subsection{Simmetria di asse un sottospazio}
\begin{definizione}
Su $\R$, la \emph{simmetria di asse $U$ e direzione $W$} è
\[
 \sigma_U^W:V\longrightarrow V,\qquad \sigma_U^W(u+w)=u-w.
\]
\end{definizione}
\begin{schema}
\begin{tikzpicture}[scale=.85]
 \coordinate (O) at (0,0); \coordinate (U) at (3,0);
 \coordinate (W) at (.7,1.3); \coordinate (V) at (3.7,1.3);
 \coordinate (S) at (2.3,-1.3);
 \draw[asse] (-.5,0)--(4.6,0) node[right] {$U$};
 \draw[asse] (-.9,-1.67)--(1.05,1.95) node[above] {$W$};
 \draw[dashed,black!45] (W)--(V)--(S)--(-.7,-1.3);
 \draw[vettore] (O)--(U) node[below right] {$u$};
 \draw[vettore] (O)--(V) node[right] {$v=u+w$};
 \draw[vettore] (O)--(S) node[right] {$\sigma_U^W(v)=u-w$};
 \draw[vettore] (O)--(W) node[left] {$w$};
\end{tikzpicture}
\end{schema}
\begin{proposizione}
Posto $\sigma=\sigma_U^W$, la simmetria è un isomorfismo e soddisfa
\begin{gather*}
 \ker\sigma=\{0_V\},\qquad \operatorname{Im}\sigma=V,\qquad
 \sigma(u)=u\quad(u\in U),\\
 \sigma\circ\sigma=\operatorname{id}_V,\qquad
 U=\operatorname{Im}(\operatorname{id}_V+\sigma),\qquad
 W=\ker(\operatorname{id}_V+\sigma).
\end{gather*}
\end{proposizione}
\begin{proof}
\textbf{Linearità e inversa.} La scomposizione di una combinazione lineare dà
\[
 \sigma(\lambda(u_1+w_1)+\mu(u_2+w_2))
 =\lambda(u_1-w_1)+\mu(u_2-w_2).
\]
Inoltre $\sigma(u)=u$ e $\sigma(u-w)=u+w$, quindi $\sigma^2=\operatorname{id}_V$. Ne segue che $\sigma$ è la propria inversa: è iniettiva e suriettiva, con nucleo nullo e immagine $V$.

\textbf{Recupero dei due sottospazi.} Per $v=u+w$ vale
\[
 (\operatorname{id}_V+\sigma)(v)=(u+w)+(u-w)=2u.
\]
Questi vettori appartengono a $U$, e ogni $u\in U$ si ottiene applicando $\operatorname{id}_V+\sigma$ a $u/2$: l'immagine è esattamente $U$. La stessa uguaglianza mostra che il nucleo è dato da $u=0$, cioè da $W$.\footnote{Le proprietà della simmetria si estendono a ogni campo di caratteristica diversa da $2$. In caratteristica $2$ si ha $u-w=u+w$, quindi la simmetria così definita coincide con l'identità e le ultime due caratterizzazioni in generale non valgono.}
\end{proof}
Le due costruzioni sono collegate da $\sigma_U^W=2\pi_U^W-\operatorname{id}_V$ e, su $\R$, da $\pi_U^W=(\operatorname{id}_V+\sigma_U^W)/2$.

% !TeX root = ../main.tex
\section{Matrice associata a un'applicazione lineare}
\subsection{Le immagini di una base determinano l'applicazione}
Sia $\mathcal B=(v_1,\ldots,v_n)$ una base di $V$. Per linearità,
\[
 v=\sum_{i=1}^n\lambda_iv_i\quad\Longrightarrow\quad
 f(v)=\sum_{i=1}^n\lambda_if(v_i).
\]
Per conoscere $f$ su tutto $V$ basta quindi conoscere le immagini dei vettori di una base. Viceversa, scelti arbitrariamente $z_1,\ldots,z_n\in W$, la formula $f(\sum_i\lambda_iv_i)=\sum_i\lambda_iz_i$ definisce un'unica applicazione lineare con $f(v_i)=z_i$: è ben definita per l'unicità delle coordinate, e la linearità si verifica sommando le coordinate.

\begin{definizione}
Fissate le basi ordinate $\mathcal B=(v_1,\ldots,v_n)$ di $V$ e $\mathcal C=(w_1,\ldots,w_m)$ di $W$, scriviamo
\[
 f(v_i)=\sum_{h=1}^m a_{hi}w_h.
\]
La \emph{matrice di $f$ nelle basi $\mathcal B,\mathcal C$} è
\[
 A=M_{\mathcal B}^{\mathcal C}(f)=(a_{hi})\in M(m,n,K).
\]
La sua colonna $i$ è $[f(v_i)]_{\mathcal C}$. La base in basso si riferisce al dominio e quella in alto al codominio.\footnote{La notazione $M_{\mathcal B}^{\mathcal C}(f)$ corrisponde alla notazione manoscritta $\psi_{\mathcal B}^{\mathcal C}(f)$. Nel testo di Bertapelle e Candilera è usata anche la notazione $\alpha_{\mathcal B\mid\mathcal C}(f)$.}
\end{definizione}
La matrice dipende da entrambe le basi; ha tante colonne quanti sono i vettori della base del dominio e tante righe quanti sono quelli della base del codominio.

\begin{proposizione}[Formula delle coordinate]
\label{prop:coordinate-applicazione}
Per ogni $v\in V$,
\[
 [f(v)]_{\mathcal C}=M_{\mathcal B}^{\mathcal C}(f)[v]_{\mathcal B}.
\]
\end{proposizione}
\begin{proof}
Posto $v=\sum_i\lambda_iv_i$, la linearità e le espressioni di $f(v_i)$ danno
\[
 f(v)=\sum_i\lambda_i\sum_h a_{hi}w_h
     =\sum_h\left(\sum_i a_{hi}\lambda_i\right)w_h.
\]
Per unicità delle coordinate, se $[f(v)]_{\mathcal C}=(\mu_h)$, allora $\mu_h=\sum_i a_{hi}\lambda_i$. È esattamente la componente $h$ del prodotto $A[v]_{\mathcal B}$.
\end{proof}

\subsection{Un'applicazione in due coppie di basi}
\begin{esercizio}
Sia $f:\R^2\to\R^2$ definita da
\[
 f(x_1,x_2)=(2x_1-x_2,\ x_1+3x_2).
\]
Determinare la matrice di $f$ rispetto alle basi
\[
 \mathcal B=\bigl((1,1),(2,-1)\bigr),\qquad
 \mathcal C=\bigl((-3,1),(2,0)\bigr),
\]
e rispetto alle basi canoniche del dominio e del codominio.
\end{esercizio}
\begin{svolgimento}
\textbf{Immagini e coordinate nella base $\mathcal C$.} La formula di $f$ è lineare nelle due variabili, dunque definisce un'applicazione lineare. Le famiglie indicate sono basi: risolvendo $\lambda(1,1)+\mu(2,-1)=0$ si ottiene $\lambda=\mu=0$, e risolvendo $\alpha(-3,1)+\beta(2,0)=0$ si trova $\alpha=\beta=0$; due vettori indipendenti in $\R^2$ formano una base.

Si ha $f(1,1)=(1,4)$. Per esprimerlo come $a_{11}(-3,1)+a_{21}(2,0)$ risolviamo
\[
 \begin{cases}-3a_{11}+2a_{21}=1,\\a_{11}=4,\end{cases}
 \qquad\Longrightarrow\qquad a_{11}=4,\quad a_{21}=\frac{13}{2}.
\]
Analogamente $f(2,-1)=(5,-1)$, e
\[
 \begin{cases}-3a_{12}+2a_{22}=5,\\a_{12}=-1,\end{cases}
 \qquad\Longrightarrow\qquad a_{12}=-1,\quad a_{22}=1.
\]
Disponendo le due colonne otteniamo
\[
 M_{\mathcal B}^{\mathcal C}(f)=\begin{pmatrix}4&-1\\13/2&1\end{pmatrix}.
\]
\textbf{Basi canoniche.} Se $\mathcal E=(e_1,e_2)$ è la base canonica, le colonne sono direttamente $f(e_1)=(2,1)$ e $f(e_2)=(-1,3)$, quindi
\[
 M_{\mathcal E}^{\mathcal E}(f)=\begin{pmatrix}2&-1\\1&3\end{pmatrix}.
\]
Le due matrici rappresentano la stessa funzione, ma agiscono su colonne di coordinate riferite a basi diverse.
\end{svolgimento}

\section{Nucleo, immagine e controimmagini mediante matrici}
\subsection{Sistemi per il nucleo e per le fibre}
Posto $A=M_{\mathcal B}^{\mathcal C}(f)$, la formula delle coordinate dà
\[
 v\in\ker f\quad\Longleftrightarrow\quad A[v]_{\mathcal B}=0.
\]
Per trovare una base di $\ker f$, risolviamo il sistema omogeneo $Ax=0$, scegliamo una base dello spazio delle soluzioni in $K^n$ e ricostruiamo i vettori di $V$ mediante la base $\mathcal B$. Le colonne risolutive sono coordinate, non necessariamente i vettori del dominio stessi.

\begin{definizione}
Per $w\in W$, la \emph{controimmagine} di $w$, o \emph{fibra sopra $w$}, è
\[
 f^{-1}(w)=\{v\in V:f(v)=w\}.
\]
Qui $f^{-1}(w)$ indica un insieme, anche quando $f$ non è invertibile; si incontra anche la scrittura $f^{\leftarrow}(w)$.
\end{definizione}
Se $b=[w]_{\mathcal C}$, calcolare $f^{-1}(w)$ equivale a risolvere $Ax=b$. La fibra è vuota se $w\notin\operatorname{Im}f$. Se invece $f(v_0)=w$, allora
\[
 f^{-1}(w)=v_0+\ker f.
\]
Infatti $f(v)=w=f(v_0)$ equivale a $f(v-v_0)=0$. È la versione astratta della struttura delle soluzioni di un sistema non omogeneo già dimostrata nel capitolo precedente.

\subsection{L'immagine e il rango della matrice}
\begin{proposizione}
Se $\mathcal B=(v_1,\ldots,v_n)$ è una base di $V$, allora
\[
 \operatorname{Im}f=\operatorname{span}\{f(v_1),\ldots,f(v_n)\}.
\]
Per ogni scelta delle basi del dominio e del codominio,
\[
 \dim\operatorname{Im}f=\operatorname{rg}M_{\mathcal B}^{\mathcal C}(f).
\]
\end{proposizione}
\begin{proof}
Ogni $f(v)$ è combinazione delle immagini dei vettori di base; viceversa ogni combinazione $\sum_i\lambda_if(v_i)=f(\sum_i\lambda_iv_i)$ appartiene all'immagine. Le colonne della matrice sono le coordinate di queste immagini. L'isomorfismo delle coordinate conserva le relazioni lineari, dunque la dimensione dello spazio generato dalle immagini è il rango per colonne della matrice, che coincide con il rango per righe.
\end{proof}
Le immagini di una base generano sempre l'immagine, ma possono essere dipendenti. Per ricavarne una base si selezionano le colonne pivot nella matrice ridotta e si prendono le immagini $f(v_i)$ corrispondenti alle colonne \emph{originarie}. Inoltre
\[
 \dim\ker f=n-\operatorname{rg}A,
\]
in accordo sia con il teorema della dimensione sia con il conteggio dei parametri di $Ax=0$.

\section{Operazioni sulle applicazioni e sulle matrici}
\subsection{Lo spazio delle applicazioni lineari}
\begin{definizione}
Si indica con $\operatorname{Hom}_K(V,W)$ l'insieme delle applicazioni lineari da $V$ a $W$. Le operazioni sono definite punto per punto:
\[
 (f+g)(v)=f(v)+g(v),\qquad (\lambda f)(v)=\lambda f(v).
\]
\end{definizione}
Queste operazioni rendono $\operatorname{Hom}_K(V,W)$ uno spazio vettoriale. Infatti, per $a,b\in K$,
\begin{align*}
 (f+g)(av+bv')&=a(f+g)(v)+b(f+g)(v'),\\
 (\lambda f)(av+bv')&=a(\lambda f)(v)+b(\lambda f)(v').
\end{align*}
La funzione nulla è lineare, l'opposto è $-f$ e gli assiomi dello spazio vettoriale si verificano valutando in ciascun $v$ e usando quelli di $W$.

\begin{proposizione}
A basi fissate, la somma e la moltiplicazione per scalari corrispondono alle stesse operazioni fra matrici:
\[
 M_{\mathcal B}^{\mathcal C}(f+g)=M_{\mathcal B}^{\mathcal C}(f)+M_{\mathcal B}^{\mathcal C}(g),\qquad
 M_{\mathcal B}^{\mathcal C}(\lambda f)=\lambda M_{\mathcal B}^{\mathcal C}(f).
\]
\end{proposizione}
\begin{proof}
La colonna $i$ della matrice di $f+g$ è $[f(v_i)+g(v_i)]_{\mathcal C}=[f(v_i)]_{\mathcal C}+[g(v_i)]_{\mathcal C}$. Per gli scalari si usa allo stesso modo $[\lambda f(v_i)]_{\mathcal C}=\lambda[f(v_i)]_{\mathcal C}$.
\end{proof}
La corrispondenza $f\mapsto M_{\mathcal B}^{\mathcal C}(f)$ è quindi un isomorfismo fra $\operatorname{Hom}_K(V,W)$ e $M(m,n,K)$: ogni matrice prescrive liberamente le immagini dei vettori di base e determina un'unica applicazione. In particolare, in dimensione finita, $\dim\operatorname{Hom}_K(V,W)=mn$.

\subsection{La composizione corrisponde al prodotto}
\begin{proposizione}
Per $f:V\to W$, $g:W\to Z$ lineari e basi $\mathcal B,\mathcal C,\mathcal D$ rispettivamente di $V,W,Z$, si ha
\[
 M_{\mathcal B}^{\mathcal D}(g\circ f)
 =M_{\mathcal C}^{\mathcal D}(g)\,M_{\mathcal B}^{\mathcal C}(f).
\]
\end{proposizione}
\begin{proof}
Per ogni $v\in V$, applicando due volte la formula delle coordinate,
\begin{align*}
 [(g\circ f)(v)]_{\mathcal D}
 &=M_{\mathcal C}^{\mathcal D}(g)[f(v)]_{\mathcal C}\\
 &=M_{\mathcal C}^{\mathcal D}(g)M_{\mathcal B}^{\mathcal C}(f)[v]_{\mathcal B}.
\end{align*}
Valutando sui vettori $v_i$ della base, le colonne dei due membri coincidono, e quindi coincidono le matrici.
\end{proof}
L'applicazione che agisce per prima è il fattore più a destra. La base $\mathcal C$ deve essere la stessa nel codominio di $f$ e nel dominio di $g$. In particolare, se $f$ è un isomorfismo, dalla composizione con $f^{-1}$ si ottiene
\[
 M_{\mathcal C}^{\mathcal B}(f^{-1})=\bigl(M_{\mathcal B}^{\mathcal C}(f)\bigr)^{-1}.
\]

% !TeX root = ../main.tex
\section{Cambiamento di base}
\subsection{La matrice dell'identità in basi diverse}
Siano $\mathcal B=(v_1,\ldots,v_n)$ e $\mathcal B'=(v'_1,\ldots,v'_n)$ due basi dello stesso spazio $V$. L'identità $\operatorname{id}_V:v\mapsto v$ non cambia il vettore, ma la sua matrice in queste due basi cambia la colonna delle coordinate.
\begin{definizione}
La \emph{matrice di cambiamento di base da $\mathcal B$ a $\mathcal B'$} è
\[
 P=M_{\mathcal B}^{\mathcal B'}(\operatorname{id}_V).
\]
La colonna $i$ di $P$ è $[v_i]_{\mathcal B'}$, e per ogni $v\in V$ vale
\[
 [v]_{\mathcal B'}=P[v]_{\mathcal B}.
\]
\end{definizione}
Il verso è precisato dall'uguaglianza: $P$ riceve coordinate nella base $\mathcal B$ e restituisce coordinate nella base $\mathcal B'$. Ponendo $Q=M_{\mathcal B'}^{\mathcal B}(\operatorname{id}_V)$ si ha invece $[v]_{\mathcal B}=Q[v]_{\mathcal B'}$. Componendo i due cambiamenti e usando la matrice dell'identità nella stessa base otteniamo
\[
 QP=I_n,\qquad PQ=I_n,\qquad
 M_{\mathcal B'}^{\mathcal B}(\operatorname{id}_V)=P^{-1}.
\]

\begin{schema}
\begin{tikzpicture}[scale=.9]
 \coordinate (O) at (0,0); \coordinate (V) at (3,2);
 \coordinate (B) at (3,0); \coordinate (C) at (0,2);
 \coordinate (Bp) at (3.6,1.2); \coordinate (Cp) at (-.6,.8);
 \draw[asse] (-.5,0)--(4.2,0) node[right] {$v_1$};
 \draw[asse] (0,-.4)--(0,2.8) node[above] {$v_2$};
 \draw[asse,draw=coloreschema] (-.6,-.2)--(4.2,1.4) node[right] {$v'_1$};
 \draw[asse,draw=coloreschema] (.3,-.4)--(-1.8,2.4) node[above] {$v'_2$};
 \draw[dashed,black!50] (B)--(V)--(C);
 \draw[dashed,coloreschema] (Bp)--(V)--(Cp);
 \draw[vettore] (O)--(V) node[above right] {$v$};
 \node[below] at (B) {$\lambda_1v_1$};
 \node[left] at (C) {$\lambda_2v_2$};
 \node[below right] at (Bp) {$\lambda'_1v'_1$};
 \node[left] at (Cp) {$\lambda'_2v'_2$};
\end{tikzpicture}

Lo stesso vettore $v$ ha due scomposizioni: $v=\lambda_1v_1+\lambda_2v_2=\lambda'_1v'_1+\lambda'_2v'_2$.
\end{schema}

\subsection{Come cambia la matrice di un'applicazione}
\begin{teorema}[Formula di cambiamento di basi]
Sia $f:V\to W$ lineare e siano $\mathcal B,\mathcal B'$ due basi di $V$ e $\mathcal C,\mathcal C'$ due basi di $W$. Poniamo
\[
 A=M_{\mathcal B}^{\mathcal C}(f),\quad
 A'=M_{\mathcal B'}^{\mathcal C'}(f),\quad
 P=M_{\mathcal B}^{\mathcal B'}(\operatorname{id}_V),\quad
 Q=M_{\mathcal C}^{\mathcal C'}(\operatorname{id}_W).
\]
Allora
\[
 A'P=QA,\qquad A'=QAP^{-1}.
\]
\end{teorema}
\begin{proof}
Siano $x=[v]_{\mathcal B}$, $x'=[v]_{\mathcal B'}$, $y=[f(v)]_{\mathcal C}$ e $y'=[f(v)]_{\mathcal C'}$. Le formule dei cambiamenti di coordinate e della matrice associata danno
\[
 x'=Px,\qquad y'=Qy,\qquad y=Ax,\qquad y'=A'x'.
\]
Calcolando $y'$ nei due modi si ottiene $A'Px=QAx$ per ogni $x\in K^n$. Le due matrici sono uguali perché agiscono allo stesso modo su tutte le colonne della base canonica; quindi $A'P=QA$. Moltiplicando a destra per $P^{-1}$ si conclude $A'=QAP^{-1}$.
\end{proof}
Esplicitando le basi, l'ultima formula diventa
\[
 M_{\mathcal B'}^{\mathcal C'}(f)
 =M_{\mathcal C}^{\mathcal C'}(\operatorname{id}_W)\,
 M_{\mathcal B}^{\mathcal C}(f)\,
 M_{\mathcal B'}^{\mathcal B}(\operatorname{id}_V).
\]
Si leggono i tre passi da destra: dalle nuove coordinate del dominio alle vecchie; applicazione di $f$ nelle vecchie basi; dalle vecchie coordinate del codominio alle nuove.

\begin{esempio}
Riprendiamo $f(x_1,x_2)=(2x_1-x_2,x_1+3x_2)$ e le basi $\mathcal B,\mathcal C$ dell'esercizio precedente. Le matrici che convertono le loro coordinate in coordinate canoniche sono
\[
 S_{\mathcal B}=\begin{pmatrix}1&2\\1&-1\end{pmatrix},\qquad
 S_{\mathcal C}=\begin{pmatrix}-3&2\\1&0\end{pmatrix}.
\]
Le colonne sono i vettori delle rispettive basi. Se $F=M_{\mathcal E}^{\mathcal E}(f)$, allora $S_{\mathcal C}M_{\mathcal B}^{\mathcal C}(f)=FS_{\mathcal B}$, per cui
\begin{align*}
 M_{\mathcal B}^{\mathcal C}(f)
 &=S_{\mathcal C}^{-1}FS_{\mathcal B}\\
 &=\begin{pmatrix}0&1\\1/2&3/2\end{pmatrix}
   \begin{pmatrix}2&-1\\1&3\end{pmatrix}
   \begin{pmatrix}1&2\\1&-1\end{pmatrix}
 =\begin{pmatrix}4&-1\\13/2&1\end{pmatrix}.
\end{align*}
Questo conferma il calcolo effettuato direttamente sulle immagini dei due vettori di base.
\end{esempio}

